\documentclass[10pt,letterpaper]{article}
\usepackage[margin=0.7in]{geometry}
\usepackage{parskip}
\usepackage{microtype}
\usepackage[T1]{fontenc}
\usepackage{lmodern}
\usepackage[compact]{titlesec}
\usepackage[hidelinks]{hyperref}
\titlespacing*{\section}{0pt}{6pt}{2pt}
\setlength{\parskip}{3pt}

\title{\vspace{-2em}\textbf{Assignment 1: Single-Server Key-Value Store}\\[1pt]
\normalsize CSCI 5144 -- Cloud Computing \quad$\cdot$\quad \url{https://github.com/Laeeq14/CSCI_5144_Assignments/}}
\author{
  \begin{tabular}{c c c}
    Laeeq Ahmed & Adam Meziou & Dang Nguyen \\
    \texttt{fazil017@umn.edu} & \texttt{mezio004@umn.edu} & \texttt{nguy6082@umn.edu} \\[6pt]
    \multicolumn{3}{c}{
      \begin{tabular}{c c}
        Yassir Abdalla & Younus Ali \\
        \texttt{abdal167@umn.edu} & \texttt{ali00380@umn.edu}
      \end{tabular}
    }
  \end{tabular}
}
\date{}

\begin{document}
\maketitle
\vspace{-2.2em}

\section*{Design Decisions}
The server is written in Go with the standard \texttt{net/http} package,
chosen for goroutine-based concurrency, built-in HTTP keep-alive, and a single
static binary with no runtime dependency.

\textbf{Concurrency.} A \texttt{map[string]string} is guarded by a
\texttt{sync.RWMutex}: GETs take a read lock and run in parallel; PUT and DEL
take the write lock, so concurrent writes to the same key are serialized.
PUT is exposed as \texttt{POST /<key>} with a JSON body, matching the
provided benchmark.

\textbf{Persistence.} The store is flushed every 10\,s and within 2\,s of any
write (fixed-deadline debounce). Each save copies the map under a read lock,
writes a temp file, calls \texttt{fsync}, and atomically renames it. On Ctrl+C
(SIGINT) a final flush runs before exit; tested by \texttt{kv\_bench.py set},
Ctrl+C, restart, \texttt{get} (\texttt{mismatch=0}).

\textbf{Logging.} Operations go through a 65,536-entry buffered channel to a
background writer (flushed every 200\,ms), keeping file I/O off the request
path. Timestamps are taken in the request goroutine. If the channel fills,
entries are dropped and counted (reported at shutdown); no drops were observed
in any benchmark run.

\textbf{Error handling.} Missing key: 404; malformed JSON or missing
\texttt{value}: 400; body over 1\,MiB: 413 (all with a JSON error body);
unsupported method: 405 (plain-text body).

\section*{Challenges}
The main challenge was separating client from server bottlenecks. The provided
benchmark uses three Python threads sharing the GIL, and on our machine
\texttt{requests} costs about 0.37\,ms of CPU per call (measured with a single
\texttt{http.client} connection against a warm server), so throughput is
client-bound at roughly 2,000--3,000\,ops/s regardless of server speed.
Running the same benchmark under WSL2 gave a similar score, which supports
the client-bound explanation rather than a Windows-specific effect.
A Go load generator reached about 169k\,ops/s at 100 goroutines
(non-verifying Go load generator on the same host, peak of the sweep); see
README.
A second issue was shutdown: waiting for the persistence goroutine with
\texttt{time.Sleep} could cut off a slow \texttt{fsync}, so it was replaced
with a done channel that \texttt{main} blocks on.

\section*{Testing and Results}
Windows\,11, Go\,1.27.0, Python\,3.14.4, laptop plugged in.
\texttt{kv\_bench.py} reported \texttt{mismatch=0, error=0} in both phases, and
data survived a Ctrl+C restart. Eight runs gave 2,062--2,221\,ops/s (median
2,101, average latency about 1.4\,ms). On battery the same benchmark gave
897--1,144\,ops/s, consistent with CPU throttling of the Python client. An
18-case edge test suite (missing keys, bad or oversized bodies, special
characters, non-ASCII values, wrong methods) passes on every run.

\section*{Assumptions}
Single process, no replication. A hard kill within the 2\,s debounce window can
lose recent writes. Bodies are limited to 1\,MiB; keys and values are UTF-8
strings. Benchmark results use the \texttt{kv\_bench.py} script provided by the
instructor with its default configuration (3 threads, 300 keys).

\section*{Potential Improvements}
A write-ahead log would remove the loss window without a full snapshot per
write; sharding the map with per-shard locks would reduce write contention;
TTLs would support cache workloads; and Raft-based replication would add fault
tolerance.

\end{document}